\section{The reduction to sparse cones and concentration}
\label{sec:reduction}

This section proves Theorem~\ref{thm:main} from the sparse-cones theorem,
the rich-motion theorem, and the uniform concentration theorem.
Their proofs occupy the later sections and appendices. We include the
selection and sampling arguments here, since keeping their constants
uniform is essential.

\subsection{The contradiction regime}

Suppose that Theorem~\ref{thm:main} fails, and take the least dimension
\(d\ge3\) in which it fails. There is then a sequence of finite sets
\(P\subset\mathbb R^d\) for which, with
\begin{equation}\label{red:regime}
 N=|P|,\qquad M=1+|\Delta(P)|,\qquad
 B=N^{1/d},\qquad A=B^2,
 \qquad N\longrightarrow\infty,\quad M/A\longrightarrow0.
\end{equation}
Here and below a statement along the sequence is understood after
discarding finitely many terms, if necessary. The fact that
\(N\to\infty\) follows because a set of at least two points determines
at least one distance. Adding \(1\) to the number of distances does
not affect the vanishing ratio.

For \(3\le u<d\), the assertion in dimension \(u\) is available by
minimality of \(d\). We also use the planar result proved in
Theorem~\ref{ci:planar-distances}. These give the following bounds.

\begin{lemma}[Occupancy of proper flats]\label{red:flat-caps}
In the regime~\eqref{red:regime}, a real affine \(u\)-flat contains
at most
\[
 \begin{cases}
  1,&u=0,\\
  2M,&u=1,\\
  C M\log(2M),&u=2,\\
  C_u M^{u/2},&3\le u<d
 \end{cases}
\]
points of \(P\). A circle contains at most \(2M\) points.
In particular there is a fixed \(c_0>0\) such that every proper real
affine flat contains at most \(N B^{-c_0}\) points for all sufficiently
large terms of the sequence. The same assertion holds for a proper
complex affine flat when only its real sample points are counted.
Planes have the more precise bound \(B^{2+o(1)}\).
\end{lemma}

\begin{proof}
From a point on a line, any positive distance has at most two realizations
on that line. For a circle, use a sample point on the circle: the
intersection with a circle of prescribed positive radius about that
point has at most two points. These observations include the sample
point itself after using \(M=1+|\Delta(P)|\).

The planar bound is Corollary~\ref{ci:planar-occupancy}. The other bounds
follow by applying the lower-dimensional assertion to the points in
the flat, whose pairwise distances are among the distances of \(P\).
Since \(M\le B^2\), the planar bound is \(O(B^2\log B)\), while a proper
flat of dimension \(u\ge3\) contains \(O(B^u)\) points. Thus, for example,
\(c_0=1/2\) works after advancing along the sequence. A complex affine
flat containing real points contains their real affine span, whose
real dimension is at most its complex dimension. Applying the real
bound to that span proves the last assertion.
\end{proof}

We use complex varieties for polynomial arguments. Degree means the
degree of the projective closure; for a family of components of a fixed
dimension it means the sum of their degrees. A projective direction
from \(p\) is denoted by
\[
 \pi_p(q)=[q-p]\in\mathbb P^{d-1}.
\]
Restriction of polynomials of degree at most \(D\) to an affine set
\(T\) has rank \(H_T(D)\). In projective space we use homogeneous
polynomials of degree exactly \(D\). The degree-\(D\) envelope
\(\Env_D(T)\) is the common zero locus of all such polynomials
vanishing on \(T\). Its rank at that degree equals \(H_T(D)\).

Here is the form of the sparse-cones theorem that we will use.

\begin{quote}
In the regime~\eqref{red:regime}, one cannot assign to each \(p\in P\)
a union \(C_p\) of irreducible projective curves, each spanning a
projective space of dimension at most two, with
\[
 \deg C_p\le\delta N/A,\qquad\delta\to0,
\]
such that all but \(o(N^2)\) unordered pairs \(\{p,q\}\) satisfy
\(\pi_p(q)\in C_p\) or \(\pi_q(p)\in C_q\).
\end{quote}
This is Theorem~\ref{sc:theorem}. Both the alternative between the two
endpoints and the absence of a prescribed rate for \(\delta\) are
important.

\subsection{A symmetric selection of directions}

\begin{proposition}\label{red:directions}
After passing to a subsequence, there is an undirected graph
\(\mathcal G\) on \(P\) with at least \(cN^2\) edges such that, for every
vertex \(p\) and every irreducible projective variety
\(Y\subset\mathbb P^{d-1}\) of dimension \(a\),
\begin{equation}\label{eq:direction-count}
 \#\{q:\{p,q\}\in\mathcal G,\ \pi_p(q)\in Y\}
 \le C\,\deg(Y)A^a.
\end{equation}
The positive constants \(c,C\) are fixed along the subsequence.
\end{proposition}

\begin{proof}
First discard pairs on lines containing more than \(H\) sample points,
where \(H\) will be a sufficiently large fixed constant.
A generic planar projection preserves the relevant distinct points,
lines, and incidences. The rich-line consequence of
Theorem~\ref{ci:st} therefore bounds the number of physical lines with
at least \(t\) points by
\[
 C(N^2/t^3+N/t),\qquad t\ge2.
\]
If \(L_{\max}=o(N)\) is the maximum occupancy, a dyadic sum gives
\[
 \#\{\text{discarded pairs}\}
 \le C N^2/H+C N L_{\max}=O(N^2/H)+o(N^2).
\]

At each \(p\), independently of the other stars, greedily remove
directions on any irreducible curve spanning at most a projective plane
that currently contains more than \(K A\) times its degree many
partners. Each removal contains at least one previously unremoved
partner, so the process terminates. Record the curves chosen in that
star. Their total degree is at most \(N/(KA)\), because their charged
partner sets are disjoint. Now keep an unordered pair only when it
survived both endpoint procedures.

There must be fixed \(H,K\) and a subsequence on which a positive
fraction of pairs survives. Otherwise, a diagonal choice with
\(H,K\to\infty\) sufficiently slowly would make the surviving fraction
tend to zero, while the first, rich-line discard also has vanishing
fraction. The recorded curves would then cover all but \(o(N^2)\)
pairs from one endpoint or the other and have total degree at most
\(N/(KA)=o(N/A)\) in every star. This contradicts
Theorem~\ref{sc:theorem}. Fix such \(H,K\) and a positive-density
subsequence.

We prove~\eqref{eq:direction-count} by induction on \(a\). For \(a=0\)
it is the light-line bound. Fix a variety \(Y\) of dimension \(a\ge1\),
a center \(p\), and the selected partners counted on \(Y\). Let \(U\)
be the real linear span of their displacements from \(p\), and put
\(\ell=\dim U\). There is nothing to prove if this collection is empty.

Consider a full radius class about \(p\), meaning all points of \(P\)
at one prescribed distance from \(p\). Suppose it has \(L\ge2M\)
distinct orthogonal projections to \(U\), and choose physical points
\(z_1,\ldots,z_L\in P\) in that class representing those projections.
For each counted partner \(q\), the distances
\(|q-z_i|\), including zero when it occurs, take at most \(M\) values.
Consequently at least
\[
 L^2/M-L\ge L^2/(2M)
\]
ordered pairs \(i\ne j\) have equal distance to \(q\).
As the \(z_i\) lie in a common radius class,
\[
 |q-z_i|=|q-z_j|
 \quad\Longrightarrow\quad
 (q-p)\cdot(z_i-z_j)=0.
\]
The projection of \(z_i-z_j\) to \(U\) is nonzero. The resulting
hyperplane section is proper on \(Y\): if its linear equation vanished
on \(Y\), it would vanish on all counted displacements and hence on
\(U\), which is impossible for its nonzero normal in \(U\).
B\'ezout's theorem and the induction hypothesis bound the number of
partners on each section by \(C\deg(Y)A^{a-1}\). Double counting gives
\[
 \#\{q\text{ counted}\}\frac{L^2}{2M}
 \le C L^2\deg(Y)A^{a-1}.
\]
This is~\eqref{eq:direction-count}, since \(M\le A\).

In the remaining case, each of the at most \(M\) radius classes has
fewer than \(2M\) projections to \(U\). Thus the full sample has
\(O(M^2)\) fibers on which both the radius about \(p\) and the
projection to \(U\) are fixed.
A fiber lies on a sphere in an affine space of dimension \(d-\ell\).
If \(\ell\ge4\) and \(d-\ell\ge3\), lower-dimensional induction bounds
each fiber by \(O(M^{(d-\ell)/2})\). If \(d-\ell=2\), the circle bound
is \(O(M)\); if \(d-\ell\le1\), the fiber has at most two points.
In all cases the total is \(O(M^{d/2})\), contradicting
\(N/M^{d/2}\to\infty\). Thus this alternative is impossible when
\(\ell\ge4\).

It remains to consider \(\ell\le3\).
For \(a\ge2\), Lemma~\ref{red:flat-caps} bounds the counted partners
by \(O(A^2)\), which is sufficient. If the span is the entire space,
then \(d=3\) and the trivial bound \(N=A^{3/2}\) is sufficient instead.
For \(a=1\), either \(Y\subset\mathbb P(U_{\mathbb C})\), in which case
the greedy curve bound applies, or its intersection with that
projective subspace is proper and consists of at most \(\deg Y\)
points. The light-line bound handles the latter case. This completes
the induction. Fixed increases in the constants at the finitely many
dimensions are harmless.
\end{proof}

\begin{proposition}[Independent evaluations in every star]
\label{red:independent-stars}
A further symmetric selection retains at least \(c'N^2\) pairs and
has the following property: the direction-evaluation vectors of the
partners in each star are linearly independent for homogeneous
polynomials of degree \(D=\lfloor A\rfloor\).
In particular each selected star admits homogeneous separators of
degree \(D\).
\end{proposition}

\begin{proof}
First consider any subset \(S\) of the partners in one star, counting
physical partners even when their projective directions agree.
Form the projective envelope of those directions at degree \(D\).
Assign each direction to an irreducible component containing it, and
group these components by dimension. For a group of total degree \(a_j\)
and dimension \(j\), Proposition~\ref{red:directions} bounds the assigned
physical partners by \(C a_j A^j\). The component Hilbert lower bound
in Lemma~\ref{poly:hilbert}, applied at the defining degree \(D\), gives
\[
 H_S(D)=H_{\Env_D(S)}(D)\ge c_d a_jD^j.
\]
There are at most \(d\) groups, and \(D\asymp A\). Therefore
\begin{equation}\label{eq:star-rank}
 H_S(D)\ge c_1|S|
\end{equation}
with a fixed \(c_1>0\), uniformly for every subset \(S\).

Take a maximal undirected subset of the edges from
Proposition~\ref{red:directions} whose evaluations are independent at
both endpoints. Each omitted edge is blocked at an endpoint: its
evaluation vector lies in the span of the selected vectors there.
Charge it to one such endpoint. If the selected degree at \(p\) is
\(s_p\), all vectors charged to \(p\) together have rank at most \(s_p\).
Equation~\eqref{eq:star-rank} bounds their number by \(s_p/c_1\).
Summing over vertices bounds the number of omitted edges by
\(2/c_1\) times the number of selected edges. Positive density follows.

Independent evaluation vectors mean that the evaluation map onto the
coordinates indexed by the star is surjective. Taking the coordinate
basis in its image gives the asserted homogeneous separators.
\end{proof}

Write \(\mathcal S\) for this final selection. A selected star has at
most one partner in any projective direction, since two proportional
evaluation vectors could not be independent.

\subsection{The equal-distance graph and its pair flats}

Take \(Q=RP\), where \(R\) is a sufficiently generic real rotation,
and transport \(\mathcal S\) to \(Q\).
On the \(N^2\) vertices \((p,q)\in P\times Q\), put an edge between
\((p,q)\) and \((p',q')\) exactly when
\[
 \{p,p'\}\in\mathcal S,\qquad
 \{q,q'\}\in R\mathcal S,\qquad |p-p'|=|q-q'|.
\]
If \(s_\rho\) is the number of selected physical pairs of length
\(\rho\), the number of these edges is
\(2\sum_\rho s_\rho^2\). Cauchy--Schwarz gives
\begin{equation}\label{eq:equality-edges}
 2\sum_\rho s_\rho^2
 \ge \frac{2|\mathcal S|^2}{M}
 \ge c_2 N^4/M.
\end{equation}

For \(d=3\), let \(g(x)=Sx+h\), with \(S\in O(3)\) and
\(h\in\mathbb R^3\), range over spatial isometries, and define its set
of matched graph vertices by
\[
 V_g=\{(p,g(p)):p\in P,\ g(p)\in Q\},\qquad k_g=|V_g|.
\]
The map \(g\mapsto R^{-1}\circ g\) preserves match counts and identifies
these motions with motions from \(P\) to itself. Since \(M\le A\)
eventually, Theorem~\ref{mot:rich} shows that the motions with
\(k_g\ge C_{\mathrm{rich}}A\) form a finite family, with both
orientations included.
Delete every existing graph edge whose endpoints both belong to one
of these sets \(V_g\). The number of deleted edges is at most
\[
 \sum_{k_g\ge C_{\mathrm{rich}}A}\binom{k_g}{2}
 \le\frac12\sum_{k_g\ge C_{\mathrm{rich}}A}k_g^2
 \le\frac{K_{\mathrm{rich}}}{2}\frac{N^4}{A}
 =o(N^4/M).
\]
Overlapping sets \(V_g\) can only decrease the actual deletion count.
After this deletion, no vertex can have at least
\(C_{\mathrm{rich}}A\) remaining neighbors which, together with that
vertex, are matches of one spatial isometry. Such a motion would have
at least \(C_{\mathrm{rich}}A\) matches, and all those incident edges
would have been deleted.

In every dimension, let \(E_0\) be the edge count at the start of
pruning, and repeatedly delete vertices of degree below
\(E_0/(2N^2)\). The total number of deleted edges is less than
\(E_0/2\), so at least one edge survives.
The resulting nonempty graph \(\mathcal H\) has minimum degree
\begin{equation}\label{eq:minimum-degree}
 \delta(\mathcal H)\ge c_3N^2/M,\qquad
 \delta(\mathcal H)/A^{d-1}\longrightarrow\infty.
\end{equation}
Only subsets of the original stars have been retained, so their
independence and separators remain valid.

Identify skew forms on \(\mathbb C^{d+1}\) with matrices
\[
 \begin{pmatrix} Z&c\\-c^{\mathsf T}&0\end{pmatrix},
 \qquad Z^{\mathsf T}=-Z,\quad c\in\mathbb C^d.
\]
The pair \((p,q)\) defines an affine flat
\begin{equation}\label{eq:pair-flat}
 F_{p,q}=\{(Z,c):Zv_0+c=y_0\},\qquad
 v_0=(p+q)/2,\quad y_0=(p-q)/2 .
\end{equation}
Its dimension is \(s_0=d(d-1)/2\). Distinct pairs give distinct flats:
at \(Z=0\) one recovers \(y_0\), and equality for every skew \(Z\)
then recovers \(v_0\).

Each \(F_{p,q}\) is the complexification of the corresponding real
point-pair flat in the construction of Bardwell-Evans and Sheffer
before their generic section
\cite[Theorem~6.1]{BardwellEvansSheffer2019}.

The parameter in the split quadric of
\(\mathbb P(\mathbb C^{d+1}\oplus(\mathbb C^{d+1})^*)\) is
\begin{equation}\label{eq:global-parameter}
 [v:y]=[(v_0,1):(y_0,-v_0\cdot y_0)].
\end{equation}
Indeed, on \(F_{p,q}\),
\[
 \begin{pmatrix}Z&c\\-c^{\mathsf T}&0\end{pmatrix}
 \binom{v_0}{1}
 =\binom{Zv_0+c}{-c\cdot v_0}
 =\binom{y_0}{-v_0\cdot y_0},
\]
where the last coordinate uses \(v_0\cdot Zv_0=0\). Conversely, the
first \(d\) coordinates of this evaluation equation give
\eqref{eq:pair-flat}. Thus \(F_{p,q}\) is exactly the evaluation flat
for this parameter. This parameter satisfies \(y(v)=0\) and \(v\ne0\).
We next check every hypothesis needed for the two different versions
of Theorem~\ref{conc:theorem}.

\subsection{Global and local concentration}

\begin{lemma}[Global concentration]\label{red:global-concentration}
For the retained pair flats, an irreducible variety \(V\) of dimension
\(s_0+h\) in skew-form space contains at most
\[
 C\deg(V)A^h
\]
of the flats, for \(0\le h\le d\), where the endpoint \(h=d\)
denotes the ambient space.
\end{lemma}

\begin{proof}
In Theorem~\ref{conc:theorem}, take \(e=d+1\), so \(r=d\).
Physical separators from Lemma~\ref{poly:separators}, one in \(p\)
and one in \(q\), have a product of degree \(O(M)\). Express them
in the coordinates \(v_0,y_0\) and homogenize with the last coordinate
of \(v\), which is \(1\) on every parameter. They give positive-degree
homogeneous separators of degree \(O(M)\le O(A)\).

Suppose parameters lie in one projective maximal isotropic plane.
Their original real representatives in~\eqref{eq:global-parameter}
belong to its underlying isotropic vector space. Polarization gives,
for any two of them,
\[
 (v_0-v'_0)\cdot(y_0-y'_0)=0,\qquad
 |p-p'|^2=|q-q'|^2.
\]
The first physical coordinate \(p\) is therefore injective on this
set of parameters, because \(p=p'\) forces \(q=q'\).
There are at most \(N\le A^{d-1}\) such parameters.
The total number of parameters is at most \(N^2=A^d\).
Thus Specification (a) of Theorem~\ref{conc:theorem} applies, with one
fixed constant absorbing the separator degree.

For \(h=0\), an irreducible variety containing a same-dimensional
flat equals that flat. For \(h=d\), the total-count bound suffices.
These give the two endpoints.
\end{proof}

Fix a retained vertex \((p,q)\). On its flat, use \(Z\) as coordinates,
so \(c=y_0-Zv_0\). A neighbor \((p',q')\), with
\(u=p'-p\) and \(w=q'-q\), cuts out
\begin{equation}\label{eq:neighbor-flat}
 Z(u+w)=u-w.
\end{equation}
Here \(u,w\) are nonzero real vectors with \(|u|=|w|\).
The rotation \(R\) may be chosen to ensure \(u+w\ne0\) simultaneously
for the finite set of relevant displacement pairs.

\begin{lemma}[Local concentration]\label{red:local-concentration}
Every nonzero polynomial of degree at most \(A\) on \(F_{p,q}\)
contains in its zero set at most \(C A^{d-1}\) of the neighbor flats
in~\eqref{eq:neighbor-flat}.
\end{lemma}

\begin{proof}
The local skew-form model has \(e=d\), \(r=d-1\), and parameters
\[
 [v:y]=[u+w:u-w].
\]
They lie on the split quadric, with nonzero \(v\).
They are distinct. Indeed, proportional parameters give
\((u',w')=\lambda(u,w)\). Since the vectors are real and nonzero,
\(\lambda\) is real; uniqueness of a selected partner in its projective
direction first gives \(u'=u\), then \(\lambda=1\) and \(w'=w\).

Multiply a homogeneous direction separator for \(u\) by one for \(w\)
and substitute \(u=(v+y)/2\), \(w=(v-y)/2\).
The product has degree \(O(A)\) and separates the corresponding
neighbor parameter from all others. No physical scale ambiguity
remains, because each selected star contains at most one physical
partner in each projective direction.

For parameters in a common maximal isotropic plane, polarization on
the original real representatives gives
\[
 (u_i-w_i)\cdot(u_j+w_j)
 +(u_j-w_j)\cdot(u_i+w_i)=0,
\]
or equivalently
\[
 u_i\cdot u_j=w_i\cdot w_j.
\]
The rule \(u_i\mapsto w_i\) is a well-defined real isometry on their
spans: every linear relation is preserved by equality of the Gram
matrices and positive definiteness of the real norm.
Extend it to some \(S\in O(d)\). Every neighbor in this isotropic
plane, together with the center, is a match for
\(x\mapsto q+S(x-p)\).

For \(d\ge4\), the number of such neighbors is at most
\(N=A^{d/2}\le A^{d-2}\). For \(d=3\), a collection with at least
\(C_{\mathrm{rich}}A\) neighbors would, together with the center,
violate the survivor property established when the rich match edges
were deleted. The surviving cap is therefore
\(O(A)\). This establishes precisely the isotropic hypothesis
\(O(A^{r-1})\) in every dimension. In particular no restriction to
orientation-preserving motions was made.

Factor the polynomial on \(F_{p,q}\) into irreducible hypersurfaces.
Each component has excess dimension \(h=d-2=r-1\) above a neighbor
flat, and lies in a hypersurface of degree at most \(A\).
Specification (b) of Theorem~\ref{conc:theorem} applies to it; that
specification has no total-count hypothesis. Summing its bound
\(C\deg(V)A^{d-2}\) over components of total degree at most \(A\)
proves the result.
\end{proof}

\subsection{One random sample completes the proof}

\begin{proposition}\label{red:main-contradiction}
The graph \(\mathcal H\) cannot satisfy both bounds in
Lemmas~\ref{red:global-concentration} and
\ref{red:local-concentration}.
\end{proposition}

\begin{proof}
Let \(\mathcal F\) be the family of retained global pair flats and
\(m=|\mathcal F|\). At degree \(D=\lfloor A\rfloor\), form the
envelope of their union. Each flat is contained in one of its
irreducible components. Assign it to one such component and group
components by their dimensions \(s_0+h\), \(0\le h\le d\).
If a group has total degree \(a_h\), global concentration bounds its
assigned flats by \(C a_h A^h\). The Hilbert lower bound at the
defining degree gives
\[
 H_{\bigcup\mathcal F}(D)
 \ge c a_hD^{s_0+h}.
\]
There are only \(d+1\) groups. Since \(D\asymp A\),
\begin{equation}\label{eq:global-rank-lower}
 H_{\bigcup\mathcal F}(D)\ge c_4mA^{s_0}.
\end{equation}
This includes the case in which the envelope is the entire ambient
space.

Choose each flat independently with a small fixed probability
\(\theta>0\). With probability at least \(1/2\), the number selected
is at most \(2\theta m\), by Markov's inequality.
On that event their restriction rank is at most
\[
 2\theta m\binom{D+s_0}{s_0}\le C_4\theta m A^{s_0}.
\]
Choose \(\theta\) so small that this is less than the right side
of~\eqref{eq:global-rank-lower}.

At any vertex the sampled neighbors have binomial distribution of
mean at least \(\theta\delta(\mathcal H)\). The probability of fewer
than half that mean is at most
\(\exp(-c\theta\delta(\mathcal H))\).
There are at most \(N^2\) vertices, and
\(\delta(\mathcal H)\gg A^{d-1}\), a fixed positive power of \(N\).
A union bound shows that, with probability tending to one, every
vertex has at least
\(\theta\delta(\mathcal H)/2\gg A^{d-1}\) sampled neighbors.
Both events thus hold for some choice of sample.

Its rank is strictly smaller than that of all flats. Hence a
polynomial \(f\) of degree at most \(D\) vanishes on every sampled
flat but does not vanish identically on some \(F\in\mathcal F\).
For each sampled neighbor \(F'\), it vanishes on the neighbor flat
\(F\cap F'\) inside \(F\). These neighbor flats are distinct by the
local parameter check above. There are more than \(C A^{d-1}\)
of them, contradicting Lemma~\ref{red:local-concentration} applied
to the nonzero restriction \(f|_F\).
\end{proof}

Proposition~\ref{red:main-contradiction}, together with the three
ingredients proved below, rules out~\eqref{red:regime} in the least
possible failing dimension. Thus there is no failing dimension.
Equivalently, in each fixed dimension the ratio
\(|\Delta(P)|/|P|^{2/d}\) has a positive infimum: if it did not, a
sequence with ratio tending to zero would give~\eqref{red:regime}.
This proves Theorem~\ref{thm:main}, including its constant-factor
strength.
